\section{What the Guarantees Say}\label{sec:min-guarantees}

The formal results can be read as one stack for the economic example.
If every realized leaf and merge preserves the economic oracle, the
root summary preserves the full-manifesto economic score. If the run
also preserves re-summarization, extra passes do not drift. If the
downstream preference objective depends on the document only through
its preserved economic value, training on root summaries and training
on full manifestos share a population target. And if preservation is
only approximate, sampled node-level audits turn the residual into an
explicit, finite-sample objective-gap certificate.
Figure~\ref{fig:min-theorem-deps} shows how the statements compose.

\begin{figure}[t]
    \centering
    \begin{tikzpicture}[
        every matrix/.style={
            matrix of nodes,
            ampersand replacement=\&,
            column sep=7mm, row sep=3.5mm,
            nodes={draw, rounded corners=1pt, inner sep=5pt,
                   minimum height=7.5mm, font=\small, align=center},
        },
        arr/.style={-{Latex[length=1.8mm]}, shorten >=2pt, shorten <=2pt},
    ]
        \matrix (m) {
            C1 $+$ C3 $+$ Assumption~\ref{ass:context} \&
                Preservation \& Thm.~\ref{thm:one-pass} \\
            $+\,$C2 \&
                Recompression \& Thm.~\ref{thm:multi-round} \\
            $+\,$oracle factorization \&
                Optimization \& Thm.~\ref{thm:pref-equiv} \\
            $+\,$transport, calibration, sampling \&
                Certificate \& Thms.~\ref{thm:unified-gap}, \ref{thm:e2e} \\
        };
        \foreach \r in {1,2,3,4} {
            \draw[arr] (m-\r-1.east) -- (m-\r-2.west);
            \draw[arr] (m-\r-2.east) -- (m-\r-3.west);
        }
    \end{tikzpicture}
    \caption{\textbf{Each theorem adds one ingredient to the bundle
    above it.} The four theorem statements of this section read as a
    stack: preservation uses only the realized local laws; repeated-round
    stability adds C2; preference alignment adds ``oracle factorization''
    (a measurability premise on the downstream preference objective);
    the finite-sample certificate adds transport, calibration, and
    sampling envelopes that turn a population distortion into a
    reportable bound.}
    \label{fig:min-theorem-deps}
\end{figure}

\paragraph{Root preservation.}
Fix a document partition and a full binary merge tree with leaves in
contiguous order. The root stored summary is written $Z^{(1)}(x)$;
re-summarization rounds $Z^{(R+1)}(x) = g(Z^{(R)}(x))$ iterate from
there.

\begin{theorem}[Inductive preservation]\label{thm:one-pass}
If the realized run satisfies \eqref{eq:leaf_sufficiency},
\eqref{eq:leaf_compatibility}, and Assumption~\ref{ass:context} for the
target oracle, then
\[
    \E\!\left[\metric\big(\fstar(Z^{(1)}(x)),\fstar(x)\big)\right] = 0.
\]
\end{theorem}

\begin{theorem}[Preservation under re-summarization]\label{thm:multi-round}
If the run additionally satisfies \eqref{eq:leaf_idempotence}, the
same equality holds for every round $R\ge 1$:
\[
    \E\!\left[\metric\big(\fstar(Z^{(R)}(x)),\fstar(x)\big)\right] = 0.
\]
\end{theorem}

Structurally this is induction on the tree.
Assumption~\ref{ass:context} lets an oracle-equivalent child substitute
for the raw child span inside its parent's context;
\eqref{eq:leaf_compatibility} propagates the substitution across each
merge; and \eqref{eq:leaf_idempotence} carries it through any number
of recompression passes.

\begin{corollary}[Schedule invariance]\label{cor:schedule}
For a fixed contiguous partition, any two reduction schedules that
satisfy the same realized local laws have the same expected oracle
value at the root.
\end{corollary}

In practice, left-to-right merges, balanced binary merges, and every
other schedule that respects the partition land on the same economic
reading, provided every realized call stays within the preservation
stack.

\paragraph{Preference alignment.}
\begin{theorem}[Preference-objective equivalence;
\textsc{PaperPreferenceStack}]\label{thm:pref-equiv}
Under the recompression stack, together with an \emph{oracle-factorization}
premise on the downstream objective---a measurability statement saying
the preference objective depends on the document only through its
oracle value $\fstar(x)$---and the generator-indexing premise of the
chosen application bundle, the full-document and summary preference
objectives share a population target: zero residual gives identical
argmin sets, and residual $\varepsilon$ makes every minimizer of the
summary objective $2\varepsilon$-optimal for the full objective.
\end{theorem}

The two cases correspond to the Lean statements
\texttt{paper\_preference\_stack\_same\_argmin} (exact) and
\texttt{paper\_preference\_stack\_summary\_argmin\_full\_epsilon}
(approximate); Appendix~\ref{app:min-proofs} gives the crosswalk.

The DPO and GRPO applications are premise packages that instantiate
this equivalence. Under their bundle assumptions, the preference
objective sees the document only through its preserved economic value,
so training on summaries aims at the same population target as
training on full manifestos.

\paragraph{Finite-sample certificate.}
Approximate preservation is handled by a method transport constant
$C_{\mathrm{meth}}$. Write
$\Delta_R = \E[\metric(\fstar(X),\fstar(Z^{(R)}(X)))]$
for the expected oracle distortion after $R$ recompression rounds of
Theorem~\ref{thm:multi-round}.

\begin{theorem}[Unified preference gap]\label{thm:unified-gap}
If the method-specific expected inner loss has transport constant
$C_{\mathrm{meth}}$ in oracle distance, then
\[
    |G_{\mathrm{meth}}| \le C_{\mathrm{meth}}\,\Delta_R.
\]
\end{theorem}

\begin{theorem}[Audited preference-gap certificate;
\textsc{PaperErrorStack}]\label{thm:e2e}
Let $\hat\Delta_R$ be the design-based estimate of $\Delta_R$ from a
tree-unit audit with logged positive propensities. On the event where
judge calibration, sampling error, and clipping are bounded by
$B_{\mathrm{cal}}$, $B_{\mathrm{est}}$, and $B_{\mathrm{clip}}$ in
objective units,
\[
    |G_{\mathrm{meth}}|
    \le
    C_{\mathrm{meth}}\hat{\Delta}_R + B_{\mathrm{cal}}
    + B_{\mathrm{est}} + B_{\mathrm{clip}}.
\]
The high-probability version of this inequality is the Lean statement
\texttt{paper\_error\_stack\_high\_prob}; the methodologist-facing
certificate is \texttt{PaperErrorCertificate}, with method-specific
end-to-end variants \texttt{dpo\_treepo\_end\_to\_end},
\texttt{grpo\_pl\_treepo\_end\_to\_end}, and
\texttt{grpo\_rl\_treepo\_end\_to\_end} (all in
Appendix~\ref{app:min-proofs}).
\end{theorem}

The four envelope terms carry the four real sources of slack.
$C_{\mathrm{meth}}\hat\Delta_R$ is the observed local distortion
converted into objective units. $B_{\mathrm{cal}}$ covers the gap
between the accessible judge and the trusted expert oracle.
$B_{\mathrm{est}}$ is sampling error from the audit design.
$B_{\mathrm{clip}}$ records any stabilisation bias from clipped
inverse-propensity weights. More local labels shrink
$B_{\mathrm{est}}$ directly but not $B_{\mathrm{cal}}$: calibration is
a judge-oracle question, not a sample-size question.

\paragraph{Conditional, not unconditional.}
These theorems are conditional on economic-state preservation at every
realized call. They do not assert that an arbitrary summarizer
preserves that state on political text; that claim is what the audit
in Section~\ref{sec:min-audit} estimates, and what
Section~\ref{sec:min-validation} cross-checks on two tasks whose
states are known in advance. Appendix~\ref{app:min-proofs} gives the
proof sketches.
